% The Slide documentDefinitions
%document
\documentclass[10pt]{beamer}
%theme
\usetheme{metropolis}
% packages
\usepackage{color}
\usepackage{listings}
\usepackage[ngerman]{babel}
\usepackage[utf8]{inputenc}
\usepackage{multicol}


% color definitions
\definecolor{mygreen}{rgb}{0,0.6,0}
\definecolor{mygray}{rgb}{0.5,0.5,0.5}
\definecolor{mymauve}{rgb}{0.58,0,0.82}

\lstset{
    backgroundcolor=\color{white},
    % choose the background color;
    % you must add \usepackage{color} or \usepackage{xcolor}
    basicstyle=\footnotesize\ttfamily,
    % the size of the fonts that are used for the code
    breakatwhitespace=false,
    % sets if automatic breaks should only happen at whitespace
    breaklines=true,                 % sets automatic line breaking
    captionpos=b,                    % sets the caption-position to bottom
    commentstyle=\color{mygreen},    % comment style
    % deletekeywords={...},
    % if you want to delete keywords from the given language
    extendedchars=true,
    % lets you use non-ASCII characters;
    % for 8-bits encodings only, does not work with UTF-8
    frame=single,                    % adds a frame around the code
    keepspaces=true,
    % keeps spaces in text,
    % useful for keeping indentation of code
    % (possibly needs columns=flexible)
    keywordstyle=\color{blue},       % keyword style
    % morekeywords={*,...},
    % if you want to add more keywords to the set
    numbers=left,
    % where to put the line-numbers; possible values are (none, left, right)
    numbersep=5pt,
    % how far the line-numbers are from the code
    numberstyle=\tiny\color{mygray},
    % the style that is used for the line-numbers
    rulecolor=\color{black},
    % if not set, the frame-color may be changed on line-breaks
    % within not-black text (e.g. comments (green here))
    stepnumber=1,
    % the step between two line-numbers.
    % If it's 1, each line will be numbered
    stringstyle=\color{mymauve},     % string literal style
    tabsize=4,                       % sets default tabsize to 4 spaces
    % show the filename of files included with \lstinputlisting;
    % also try caption instead of title
    language = Python,
	showspaces = false,
	showtabs = false,
	showstringspaces = false,
	escapechar = ,
}

\def\ContinueLineNumber{\lstset{firstnumber=last}}
\def\StartLineAt#1{\lstset{firstnumber=#1}}
\let\numberLineAt\StartLineAt



\newcommand{\codeline}[1]{
	\alert{\texttt{#1}}
}


% Author and Course information
% This Document contains the information about this course.

% Authors of the slides
\author{Claas de Boer, Tilman Hinnerichs}

% Name of the Course
\institute{Python-Grundlagen}

% Fancy Logo 
\titlegraphic{\hfill\includegraphics[height=1.25cm]{../images/fsr_logo_cropped}}

% Custom Bindings
% \newcommand{\codeline}[1]{
%	\alert{\texttt{#1}}
%}


% Presentation title
\title{Funktionen und Modularisierung}
\date{\today}

\begin{document}

\maketitle

\begin{frame}{Rückblick}
    \begin{itemize}
        \item Datentypen: \texttt{int, str, list, set, dict, ...}
        \item Komperatoren: \texttt{==, !=, ...}
        \item Schleifen: \texttt{while, for}
        \item Ein-/Ausgabe: \texttt{input, print}
    \end{itemize}
\end{frame}

\section{Heute: Funktionen}

\begin{frame}{Funktionen}
    \begin{itemize}
        \item Eine \textbf{Funktion} ist ein in sich geschlossener Codeblock, der eine bestimmte Aktion ausführt.
        \item \textbf{Funktionen} können mit \textbf{Parametern} aufgerufen werden und können \textbf{Rückgabewerte} liefern.
        \item \textbf{Funktionen} erlauben es Code zu \alert{modularisieren}, zu \alert{organisieren} und \alert{wiederverwendbar} zu machen.
    \end{itemize}
\end{frame}

\begin{frame}[fragile]{Built-In Funktionen}
\begin{lstlisting}
# Die Anzahl der Elemente eines Containers bestimmen
numbers = [1, 2, 3, 4, 5]
len(numbers)

# Den Typ einer Variable herausfinden
n_float = 42.8
type(n_float)

# Informationen über ein Objekt bekommen
set_of_even_numbers = set(n for n in range(1, 10) if n % 2 == 0)
help(set_of_even_numbers)

# Weitere: sorted(), enumerate(), ...
\end{lstlisting}
\textbf{\href{https://docs.python.org/3/library/functions.html}{Hier findet ihr eine Liste der Python 3.9.1 Built-in Funktionen.}}
\end{frame}

\begin{frame}[fragile]{Eigene Funktionen I}
\begin{lstlisting}
# Funktionsdefiniton startet mit dem Keyword 'def'

def simple_function():

    # Hier beginnt der Körper (body) der Funktion

    print("The mitochondria is the powerhouse of the cell.")

    # Hier endet der Körper der Funktion
\end{lstlisting}
\begin{itemize}
    \item Kein \textbf{Parameter}, kein \textbf{Rückgabewert}
\end{itemize}
\end{frame}

\begin{frame}[fragile]{Eigene Funktionen II}
\begin{lstlisting}
# Der Wert eines übergebenen Parameters wird an eine
# Variable im Körper der aufgerufenen Funktion gebunden

def format_greeting(name):
    return f"Hello {name}. Nice to meet you again!"

# Der Rückgabewert einer Funktion kann
# in einer Variable gespeichert werden

result = format_greeting("Sun")
print(result)
\end{lstlisting}
\begin{itemize}
    \item Ein \textbf{Parameter}, ein \textbf{Rückgabewert}
\end{itemize}
\end{frame}

\begin{frame}[fragile]{Eigene Funktionen III}
\begin{lstlisting}
def calculate_price(item, price, amount=5):
    costs = amount * price
    return f"{amount} {items}'s will cost you {costs} $."
\end{lstlisting}
\begin{itemize}
    \item Funktionen lassen sich auf verschiedene Arten aufrufen.
\end{itemize}
\begin{lstlisting}
>>> calculate_price("Apple", 0.3, 10)
# 10 Apple's will cost you 3.0 $

>>> calculate_price(price=0.3, amount=10, item="Apple")
# 10 Apple's will cost you 3.0 $

>>> calculate_price("Apple", 0.3)
# 10 Apple's will cost you 1.5 $
\end{lstlisting}
\end{frame}

\begin{frame}[fragile]{Parametrisieren will gekonnt sein!}
Bei den Parametern einer Funktion unterscheidet man zwischen verschiedenen Typen.
\begin{itemize}
    \item \textbf{Positionelle} Parameter (Erforderliche Parameter)
    \item \textbf{Keyword} Parameter
    \item \textbf{Optionale} Parameter (Parameter mit Default Wert)
\end{itemize}
\end{frame}

\begin{frame}[fragile]{Mini-Exkurs: Documentation}
    \begin{lstlisting}
    def complicated_funktion(gamma, offset, inertia):
        """This function computes the result
        of calculation foobar for some given
        parameters.

        A function can get really complicated!
        But we can describe what a function does in
        plain text in the function docstring.
        """
        ...
    \end{lstlisting}
    \begin{itemize}
        \item \textbf{\alert{Code wird häufiger gelesen als geschrieben.}}
        \item \textbf{Dokumentiert euren Code} und ihr müssst weniger lesen!
        \item Gebt euren Funktionen und ihren Parametern \textbf{deskriptive Namen}!
    \end{itemize}
\end{frame}

\begin{frame}{Einige Aufgaben\texttrademark 6}
    \begin{enumerate}
        \item Schreibe eine Funktion, welche die Anzahl der Vokale einer Zeichenkette berechnet.
        \item Schreibe eine Funktion, welche die Länge eines Videos im Format \texttt{mm:ss} (bspw. \texttt{02:54}) entgegen nimmt und die Länge des Videos in Sekunden zurückgibt.
        \item Schreibe eine Funktion, welche die \texttt{Fakultät} einer Ganzahl berechnet. Die Fakultät einer Ganzzahl ist das Produkt der Ganzzahl mit allen kleineren positiven Ganzzahlen.
        \item Schreibe eine Funktion, welche die \texttt{n-te} Zahl der \href{https://de.wikipedia.org/wiki/Fibonacci-Folge}{\texttt{Fibonacci-Folge}} berechnet.
    \end{enumerate}
\end{frame}

\begin{frame}{Ressourcen zum Üben/Informieren}
    \begin{itemize}
        \item \textbf{\href{https://realpython.com/defining-your-own-python-function/}{Real Python Artikel: Defining your own python function}}
        \item \textbf{\href{https://adventofcode.com/}{Advent of Code}}
    \end{itemize}
\end{frame}

\end{document}
